\scanpage{24}

% L1
\uline{Structure représentée par} $B_G=\widehat C\qquad C=\deleted{\unsure{(G,}}(\mathord{.},G)$

% L2
Relativisation dans un topos.

% L3
\uline{Topos produit}.

% L4–8
Remarques sur les $T(E)\simeq\underline{\mathrm{Hom}}\,\mathrm{top}(E,R)^o\simeq\underline{\mathrm{Hom}}_\Delta(R,E)$ : ces catégories \uline{sont stables par \uline{lim} filtrantes} (ça \unsure{t'élimine|élimine|il vient} : pas mal de conditions sur des structures : elles ne correspondent pas à des théories algébriques sur \deleted{quel que ce soit} \inserted{quelque} $\Delta'\supset\Delta$.

% L8, right
\uline{NB} En général aucun autre type de lim.

% L9–11
\uline{Question} Si $\forall E$, $T(E)$ admet des \deleted{produits fibrés quelconques} \inserted{$\underleftarrow{\lim}$ finies}, \inserted{(et les $T_R(E)\longrightarrow T_R(E')$ y commutent)} aura-t-on $T=$ restriction à $\Delta$ d'une théorie algébrique sur $\Delta_0$, i.e. $R\simeq\widehat C$, avec $C$ \inserted{(petite)} $\in\operatorname{Ob}\Delta_0$ i.e. $C$ stable par $\underleftarrow{\lim}$ finies ?

% L12–14
Si $T$ est de ce type, on sait que $T(E)$ est \unsure{$\simeq$} stable par $\underrightarrow{\lim}$ \deleted{\unsure{fini|finies}} quelconques et accessible, \deleted{(\uline{NB} Mor$(\unsure{C})$} et (donc) stable par $\underleftarrow{\lim}$ quelc.

% L15
$T_R(\mathrm{Ens})\simeq\underline{\mathrm{Pt}}(R)^o\simeq\underline{\mathrm{Fib}}(R)$. Peut être vide et $R\not\simeq\operatorname{Top}(\varnothing)$

% L16–19
\uline{Sous topos du topos libre \inserted{$R$} engendré par un objet} $X$ (espèces de structures vides sur un objet de base) : \deleted{\unsure{comme|comparer}} $R=\widehat C$, $C=$ catégorie opposée des ens finis, ou $R=$ ens. semi-simpliciaux. \inserted{\uline{Pts} : les Ens} Sous-Topos \inserted{ouverts} $=$ cribles de $\mathbf C$, on trouve les sous-topos suivants (ils sont caractérisés par leurs points) :

% L19–24, table
\diagram{$\begin{array}{r|lll}
 & \unsure{\mathbb R|\mathbf R|R}, & \varnothing, & \mathbb R^*=\unsure{\operatorname{Im}|I\mathrm m}(X\to e)\\
 \substack{\text{pts}\\\text{du p.f.}\\\text{(crible}}\ \}\to &
 \mathbb N, & \varnothing, & \mathbb N^*\\
 &\mathbf S, & \varnothing, & (\mathrm{Ens}_{f}\ \text{non vides})^*\\
 \text{Sous-théorie} & \unsure{T(E)=\mathrm{Ens}|T(E)=E|\deleted{T(E)=E}}, & \varnothing, & T(E)\subset E\\
 &&&\substack{\text{ss catégorie pleine}\\\text{des objets couvrant }E}
 \end{array}$}

% L21, left
\margin{$S=(\mathrm{Ens}_{f})^o$}

% L24–25, arrow below empty column
\diagram{\begin{tikzcd}[ampersand replacement=\&,row sep=tiny]
 {}\arrow[d] \\ \substack{T(E)=\varnothing\ \text{sauf}\\\text{si } E\simeq\mathrm{Top}(\varnothing)}
\end{tikzcd}}

% L26–28
\uline{Sous-topos fermés complémentaires}
\diagram{$\begin{array}{r|lll}
 &\varnothing,&\mathbb R,&\unsure{\mathbf R_0|\Phi_0}=\unsure{\complement|f|\mathcal C}U\\
 \text{p.p.f.}\Rightarrow &\varnothing,&\mathbb N,&\{\Delta_0\}\\
 &&&T(E)\subset E\\
 &&&\substack{\text{catég. ponctuelle}\\\text{réduite à }\deleted{\unsure{\text{l'obj}}}\text{l'objet vide}}
 \end{array}$}
% archive sheet number: 23
